\documentclass[11pt]{article}
\usepackage{amsmath, amssymb}
\usepackage{geometry}
\usepackage{enumitem}
\usepackage{hyperref}
\geometry{margin=1in}

\title{\textbf{CSE400 – Fundamentals of Probability in Computing}}
\author{\textbf{Moksha Sheth} \\\textbf{Lecture Scribe-1}}

\date{February 7 $2026$}

\begin{document}
\maketitle

\section*{Lecture 5: Bayes’ Theorem, Random Variables, and Probability Mass Function (PMF) (January 20, 2026)}\\

\hrule
\vspace{1.5cm}


\section*{1. Overview of the Lecture}

This lecture covers three connected ideas used frequently in probability for computing:

\begin{itemize}
\item \textbf{Bayes’ Theorem} (as weighted conditional probabilities + formal law of total probability + Bayes formula)
\item \textbf{Random Variables} (why we define them and how their distributions are described)
\item \textbf{Probability Mass Function (PMF)} for discrete random variables (definition + examples)
\end{itemize}\\

\vspace{1cm}
\section*{2. Key Definitions and Terminology}

\vspace{0.5cm}
\subsection*{2.1 \underline{Mutually Exclusive (Disjoint) Events}}

If events cannot happen together, they are \textbf{mutually exclusive}. In the lecture, $(AB)$ and $(AB^c)$ are disjoint.

\vspace{0.5cm}
\subsection*{2.2 \underline{Random Variable}}

A \textbf{random variable} is a real-valued function defined on the sample space; its value is determined by the outcome of the experiment.

\vspace{0.5cm}
\subsection*{2.3 \underline{Distribution of a Random Variable}}

For each possible value $(a)$, the quantity $\Pr[X=a]$ gives the probability that the random variable equals $(a)$. The lecture notes that the distribution can be visualized using a \textbf{bar diagram} where bar height at $(a)$ equals $\Pr[X=a]$.

\vspace{0.5cm}
\subsection*{2.4 \underline{Discrete Random Variable}}

A random variable is \textbf{discrete} if it can take \textbf{at most a countable} number of possible values.

\vspace{0.5cm}
\subsection*{2.5 \underline{Probability Mass Function (PMF)}}

For a discrete random variable $(X)$ with range $(R_X=\{x_1,x_2,\dots\})$, the function assigning probabilities to each possible value is called the \textbf{PMF} of $(X)$.

\vspace{0.5cm}
\subsection*{2.6 \underline{Apriori vs Posteriori Probability (as stated in lecture)}}

\begin{itemize}
\item $\Pr(B_i)$: \textbf{apriori} probability (from presupposed/self-evident model)
\item $\Pr(B_i\mid A)$: \textbf{posteriori} probability (after observing event $(A)$)
\end{itemize}

\vspace{1cm}
\section*{3. Explanation of Core Concepts and its Understanding}

\vspace{0.5cm}
\subsection*{3.1 \underline{ Bayes’ Theorem as a Weighted Average of Conditional Probabilities}}

Key idea: write $(A)$ as
\[
A = AB \cup AB^c
\]

so outcomes in $(A)$ are either “in $(A)$ and $(B)$” or “in $(A)$ and not $(B)$”.

Because $(AB)$ and $(AB^c)$ are mutually exclusive,
\[
\Pr(A)=\Pr(AB)+\Pr(AB^c).
\]

Then using conditional probability forms (as shown in the lecture), we get:
\[
\Pr(A)=\Pr(A\mid B)\Pr(B)+\Pr(A\mid B^c)\bigl(1-\Pr(B)\bigr).
\]

This is interpreted in the lecture as: \textbf{$\Pr(A)$ is a weighted average of conditional probabilities}, weighted by how likely the conditioning events are.

\vspace{0.5cm}
\subsection*{3.2 \underline{Law of Total Probability}}

The lecture introduces the law of total probability as a formal tool.

If $(B_1,\dots,B_n)$ are mutually exclusive and partition the sample space, then:
\[
\Pr(A)=\sum_{i=1}^{n}\Pr(A\mid B_i)\Pr(B_i).
\]

(Referenced in lecture as “law of total probability [Formula 3.4]”.)

\vspace{0.5cm}
\subsection*{3.3 \underline{Bayes Formula (Formal Bayes’ Theorem)}}

Using the relationship (stated in the lecture):
\[
\Pr(AB_i)=\Pr(B_i\mid A)\Pr(A),
\]

the lecture arrives at \textbf{Bayes Formula [Proposition 3.1]}.

Standard form consistent with the lecture’s setup:
\[
\Pr(B_i\mid A)=\frac{\Pr(A\mid B_i)\Pr(B_i)}{\sum_{j=1}^{n}\Pr(A\mid B_j)\Pr(B_j)}.
\]

Also recall: $\Pr(B_i)$ is apriori and $\Pr(B_i\mid A)$ is posteriori as described.

\vspace{0.5cm}
\subsection*{3.4 \underline{Random Variables: Why we define them}}

Instead of tracking the full outcome, we often track a \textbf{function} of the outcome (examples given: sum of two dice; number of heads). These functions are random variables.

\vspace{0.5cm}
\subsection*{3.5 \underline{Distribution Visualization (Bar Diagram idea)}}

\begin{itemize}
\item x-axis: values the RV can take
\item bar height at value $(a)$: $\Pr[X=a]$
\item these probabilities come from the probability of the corresponding event in the sample space
\end{itemize}

\vspace{0.5cm}
\subsection*{3.6 \underline{PMF Concept (Discrete Case)}}

A discrete random variable has countable possible values; PMF assigns probability to each single value.

\vspace{1cm}
\section*{4. Important Formulas and Mathematical Expressions}

\vspace{0.5cm}
\subsection*{4.1 \underline{Event Decomposition}}

\[
A = AB \cup AB^c
\]

\vspace{0.5cm}
\subsection*{4.2 \underline{Weighted Average Form (Two-case total probability)}}

\[
\Pr(A)=\Pr(A\mid B)\Pr(B)+\Pr(A\mid B^c)\bigl(1-\Pr(B)\bigr)
\]

\vspace{0.5cm}
\subsection*{4.3 \underline{Law of Total Probability (General partition)}}

\[
\Pr(A)=\sum_{i=1}^{n}\Pr(A\mid B_i)\Pr(B_i)
\]

\vspace{0.5cm}
\subsection*{4.4 \underline{Bayes Formula (Proposition 3.1)}}

\[
\Pr(B_i\mid A)=\frac{\Pr(A\mid B_i)\Pr(B_i)}{\sum_{j=1}^{n}\Pr(A\mid B_j)\Pr(B_j)}
\]

\vspace{0.5cm}
\subsection*{4.5 \underline{Conditional Probability used in examples}}

\[
\Pr(A\mid A_1)=\frac{\Pr(AA_1)}{\Pr(A_1)}
\]

\vspace{0.5cm}
\subsection*{4.6 \underline{Random Variable probabilities (distribution statement used repeatedly)}}

\[
\Pr[X=a]
\]

(used for bar diagram interpretation)

\vspace{0.5cm}
\subsection*{4.7 \underline{PMF example family (given)}}

\[
p(i)=c\frac{\lambda^i}{i!},\quad i=0,1,2,\dots,\ \lambda>0
\]

\vspace{1cm}
\section*{5. Derivations / Proofs (If Present in Lecture)}

\vspace{0.5cm}
\subsection*{5.1 \underline{Derivation: “Weighted Average of Conditional Probabilities”}}

\textbf{Start:}
\[
A = AB \cup AB^c
\]

\textbf{Step 1 (mutual exclusivity):}
\[
\Pr(A)=\Pr(AB)+\Pr(AB^c)
\]

\textbf{Step 2 (convert to conditional probabilities):}
\[
\Pr(AB)=\Pr(A\mid B)\Pr(B),\qquad \Pr(AB^c)=\Pr(A\mid B^c)\Pr(B^c)
\]

and $\Pr(B^c)=1-\Pr(B)$. This yields the lecture’s final form:
\[
\Pr(A)=\Pr(A\mid B)\Pr(B)+\Pr(A\mid B^c)\bigl(1-\Pr(B)\bigr).
\]

\vspace{0.5cm}
\subsection*{5.2 \underline{Derivation: Determining the constant $(c)$ in $\left(p(i)=c\frac{\lambda^i}{i!}\right)$}}

\textbf{Given PMF form:} $\left(p(i)=c\frac{\lambda^i}{i!}\right)$, $(i=0,1,2,\dots)$.

\textbf{Start with PMF property:}
\[
\sum_{i=0}^{\infty}p(i)=1
\]

So,
\[
\sum_{i=0}^{\infty}c\frac{\lambda^i}{i!}=1
\quad\Rightarrow\quad
c\sum_{i=0}^{\infty}\frac{\lambda^i}{i!}=1.
\]

\textbf{Use the series identity (as used in the lecture’s worked slide):}
\[
e^\lambda=\sum_{i=0}^{\infty}\frac{\lambda^i}{i!}
\]

Therefore,
\[
c e^\lambda=1
\quad\Rightarrow\quad
c=e^{-\lambda}.
\]

(This matches the final constant shown in the lecture’s worked example slide.)

\vspace{1cm}
\section*{6. Worked Examples / Applications}

\vspace{0.5cm}
\subsection*{6.1 \underline{Example 3.1 (Part 1/2): Insurance (Total Probability idea)}}

\textbf{Problem (given):} Accident-prone class vs not accident-prone class;

\begin{itemize}
\item $\Pr(\text{accident within 1 year} \mid \text{accident prone})=0.4$
\item $\Pr(\text{accident within 1 year} \mid \text{not accident prone})=0.2$
\item $\Pr(\text{accident prone})=0.3$
Find $\Pr(\text{accident within 1 year})$.
\end{itemize}

\textbf{Solution (as in lecture):} Let $(A_1)$ = accident within a year, $(A)$ = accident prone.
\[
\Pr(A_1)=\Pr(A_1\mid A)\Pr(A)+\Pr(A_1\mid A^c)\Pr(A^c)
\]
\[
=(0.4)(0.3)+(0.2)(0.7)=0.26.
\]

\textbf{Final answer:} $\boxed{0.26}$.

\vspace{0.5cm}
\subsection*{6.2 \underline{Example 3.1 (Part 2/2): Insurance (Bayes / Posterior probability)}}

\textbf{Problem (given):} Given that the policyholder had an accident within a year, find the probability they are accident prone.

\textbf{Solution (as in lecture):}
\[
\Pr(A\mid A_1)=\frac{\Pr(AA_1)}{\Pr(A_1)}
=\frac{\Pr(A)\Pr(A_1\mid A)}{\Pr(A_1)}
=\frac{(0.3)(0.4)}{0.26}=\frac{6}{13}.
\]

\textbf{Final answer:} $\boxed{\Pr(A\mid A_1)=6/13}$.

\vspace{0.5cm}
\subsection*{6.3 \underline{Example 3.2: Three cards (Bayes reasoning with “equally likely outcomes”)}}

\textbf{Problem (given):} Three cards: RR, BB, RB. Pick one at random, place it down. If the upper side is red, find probability the other side is black.

\textbf{Setup (as in lecture):} Let events be RR, BB, RB and let $(R)$ be “upturned side is red”. The lecture sets up:
\[
\Pr(RB\mid R)=\frac{\Pr(R\mid RB)\Pr(RB)}{\Pr(R\mid RR)\Pr(RR)+\Pr(R\mid RB)\Pr(RB)+\Pr(R\mid BB)\Pr(BB)}.
\]

With:
\begin{itemize}
\item $\Pr(R\mid RB)=1/2$, $\Pr(RB)=1/3$
\item $\Pr(R\mid RR)=1$, $\Pr(RR)=1/3$
\item $\Pr(R\mid BB)=0$, $\Pr(BB)=1/3$
\end{itemize}

\textbf{Final result (lecture conclusion):}

The probability equals $(1/3)$. The lecture explains the common wrong guess $(1/2)$ happens by incorrectly assuming the remaining possibilities are equally likely; correct reasoning counts equally likely “sides” outcomes and shows only one favorable outcome $(R3)$ among three red-up outcomes $(R1, R2, R3)$.

\textbf{Final answer:} $\boxed{1/3}$.

\vspace{0.5cm}
\subsection*{6.4 \underline{Random Variable Example 1: Tossing 3 fair coins}}

\textbf{Problem (given):} Toss 3 fair coins. Let $(Y)$ be the number of heads.

\textbf{Values and probabilities (as listed):}
\begin{itemize}
\item $\Pr(Y=0)=1/8$
\item $\Pr(Y=1)=3/8$
\item $\Pr(Y=2)=3/8$
\item $\Pr(Y=3)=1/8$
\end{itemize}

Also, since $(Y)$ must take one of these values, total probability sums to 1 (stated in the lecture).

\vspace{0.5cm}
\subsection*{6.5 \underline{PMF Example Mentioned: Two Independent Tosses of a Fair Coin}}

The lecture includes a slide titled \textbf{“Example: Two Independent Tosses of a Fair Coin”} under PMF.

(Only the example heading is present in the lecture PDF text at that point.)

\vspace{0.5cm}
\subsection*{6.6 \underline{PMF Example: $\left(p(i)=c\frac{\lambda^i}{i!}\right)$ — find $\Pr(X=0)$ and $\Pr(X>2)$}}

\textbf{Given:} $\left(p(i)=c\frac{\lambda^i}{i!}\right)$, $(i=0,1,2,\dots)$, $(\lambda>0)$.

\textbf{Step 1: Find $(c)$} (from derivation above): $(c=e^{-\lambda})$.

\textbf{Step 2: Compute $\Pr(X=0)$}
\[
\Pr(X=0)=p(0)=c\frac{\lambda^0}{0!}=c=e^{-\lambda}.
\]

\textbf{Step 3: Compute $\Pr(X>2)$}

Use complement form (as shown conceptually in the worked slide):
\[
\Pr(X>2)=1-\Pr(X\le 2)=1-\bigl(\Pr(X=0)+\Pr(X=1)+\Pr(X=2)\bigr).
\]

Where:
\[
\Pr(X=1)=p(1)=e^{-\lambda}\frac{\lambda}{1!}=e^{-\lambda}\lambda,\quad
\Pr(X=2)=p(2)=e^{-\lambda}\frac{\lambda^2}{2!}=e^{-\lambda}\frac{\lambda^2}{2}.
\]

So,
\[
\boxed{\Pr(X>2)=1-e^{-\lambda}\left(1+\lambda+\frac{\lambda^2}{2}\right)}.
\]

\vspace{1cm}
\section*{7. Important Observations and Results}

\begin{itemize}
\item \textbf{Weighted-average interpretation:} $\Pr(A)$ can be expressed as a weighted average of $\Pr(A\mid B)$ and $\Pr(A\mid B^c)$, with weights $\Pr(B)$ and $(1-\Pr(B))$.
\item \textbf{Common Bayes pitfall (cards example):} After seeing “red”, the remaining cases are not automatically equally likely; you must account for the correct underlying equally-likely outcomes.
\item \textbf{Random variables shift focus:} we track a numerical quantity derived from outcomes (sum, count, etc.), not necessarily the full outcome.
\item \textbf{Distribution view:} A distribution can be represented via bars: value vs $\Pr[X=a]$.
\item \textbf{Discrete RV $\rightarrow$ PMF:} Discrete means countable outcomes, enabling PMF assignment to each value.
\end{itemize}

\vspace{1cm}
\section*{8. Quick Revision Points}

\begin{itemize}
\item \textbf{Decompose event:} $(A=AB\cup AB^c)$.
\item \textbf{Two-case total probability:}  
$\Pr(A)=\Pr(A\mid B)\Pr(B)+\Pr(A\mid B^c)(1-\Pr(B))$.
\item \textbf{Law of total probability (partition):}  
$\Pr(A)=\sum_i \Pr(A\mid B_i)\Pr(B_i)$.
\item \textbf{Bayes formula:}  
$\Pr(B_i\mid A)=\dfrac{\Pr(A\mid B_i)\Pr(B_i)}{\sum_j \Pr(A\mid B_j)\Pr(B_j)}$.
\item \textbf{Insurance example results:}  
$\Pr(A_1)=0.26$, and $\Pr(A\mid A_1)=6/13$.
\item \textbf{Cards example result:} $\Pr(\text{RB}\mid R)=1/3$.
\item \textbf{3-coin RV distribution:} $(1/8, 3/8, 3/8, 1/8)$ for $(Y=0,1,2,3)$.
\item \textbf{PMF family:} If $\left(p(i)=c\frac{\lambda^i}{i!}\right)$, then $(c=e^{-\lambda})$, $\Pr(X=0)=e^{-\lambda}$, and
$\Pr(X>2)=1-e^{-\lambda}\left(1+\lambda+\frac{\lambda^2}{2}\right)$.
\end{itemize}

\end{document}
